\documentclass[11pt,a4paper]{article}
\usepackage[margin=27mm]{geometry}
\usepackage{amsmath,amssymb,amsthm,booktabs,longtable,hyperref}
\newtheorem{lemma}{Lemma}\newcommand{\Sz}{\operatorname{Sz}}
\newcommand{\F}{\mathbb F}
\title{Rank-three Suzuki geometries: results, corrections and realizations}
\author{Human Compression project}\date{}
\begin{document}\maketitle
\begin{abstract}
We give a complete disposition of the mathematical assertions in the fixed
1999 rank-three Suzuki article. The prime-degree obstruction and the
three necessary patterns for prime split-torus order are proved. The
composite-degree list and each of the general local-transitivity
conclusions require corrections. We supply the individual realization
arguments, finite counterexamples to thinness in both prime branches,
and a uniform scalar-subspace construction of unbounded algebraic local
rank. The classical Suzuki structure and the earlier results explicitly
quoted by the source remain attributed premises. This is a treatment of
one fixed original, not a complete replacement classification of all
Suzuki incidence geometries.
\end{abstract}

\section{Source, conventions and proof organization}
The original is D.~Leemans, \emph{The rank 3 geometries of the simple
Suzuki groups $\Sz(q)$}, Note di Matematica 19(1) (1999), 43--63.
Its twenty-one physical pages correspond to printed pages 43--63.
The principal failure of its exhaustive Theorem~5.5 has already been
established by the compatible-tori counterexample; we retain that proof.
The present treatment completes the surviving assertions of the same
original and does not count a second principal refutation.

The source uses labelled left cosets $gG_i$ with incidence defined by
nonempty intersection. On pp.~45--46 it defines the residue of a flag
inside the distinguished chamber by intersections of the parabolics.
Residual weak primitivity means that for every proper type subset $J$
some unused type gives a maximal proper inclusion
$G_{J\cup\{i\}}<G_J$. Local double transitivity concerns the resulting
algebraic rank-one residues. Actual thinness, in contrast, counts all
chambers over an incidence edge. Equivalence with definitions using
full incidence residues is qualified by flag transitivity in the source.
We preserve that distinction throughout.

The accompanying volume contains the full proof texts. Three earlier
notes on elementary Suzuki structure, normalizer families and subfield
families are reused as common theory, together with the earlier
scalar-subspace and standard-subfield arguments. Their classical inputs
are the Suzuki matrix and ovoid model, torus partition, subgroup
classification, conjugacy and automorphism data. The present work does
not reprove those classical classification theorems. It does prove the
maximality, intersection, scalar-space and field-coordinate consequences
used here. The numerical tables and twenty-nine incidence diagrams
belonging to the separate original L115 are not assumptions in these
parametric arguments and receive no completion credit here.

\section{The surviving restriction theorems}
The strict-intersection argument applies directly to the source axiom:
if $G_J\le G_i$ for an unused type $i$, then
$G_{I\setminus\{i\}}=G_I$, contradicting maximal proper inclusion at
$I\setminus\{i\}$. The point-stabilizer argument then excludes a full
point stabilizer as a parabolic in every rank at least two.

The prime-degree note proves Theorem~5.4 in its native full range:
if $q=2^n$, $n$ is prime and $q-1$ is composite, there is no such
pregeometry of rank at least three. It does not merely exclude ranks
greater than three. The prime-pattern note proves all three necessary
patterns in Theorem~5.3, the rank bound three when $q-1$ is prime,
and the rank-two local-transitivity counts: one class when $q>8$ and
three when $q=8$, in the source's specified-group equivalence convention.
A sparse matrix entry supplies the missing geometric exclusion in the
necessity argument; sufficient examples are not used as a classification.

The general normalizer note proves the content of Lemmas~4.2--4.6
and Lemma~5.1. Its cyclic and odd-dihedral induction supplies directly
the maximality consequence quoted from the earlier dihedral catalogue.
Two compatible maximal subfields and an elementary size inequality
exclude a full nonsplit normalizer. The full split normalizer can occur
only up to rank three, and its composite-parameter rank-three case has
the stated subfield/affine pattern. Consequently every maximal ambient
parabolic in rank at least four is a simple subfield group.

There is a necessary qualification to Lemma~4.5. Its printed phrase
``maximal rank ... is three'' is an upper bound, attained for example
at $q=8$. It is not an assertion that every fixed $q$ attains rank
three: Theorem~5.4 excludes attainment in its prime-degree,
composite-$q-1$ case.

\section{The quoted rank-two tables and preliminary lemmas}
Table~1 is classical maximal-subgroup data, attributed by the source
to Suzuki. The source's Theorem~5.1 quotes the earlier dihedral paper
[20], catalogue L109, which is outside the fixed accessible corpus.
Its full catalogue and diagrams remain an explicit cited premise.
The elementary cyclic and odd-dihedral results needed by the present
proofs are independently established in our normalizer note.
Likewise the observation after Theorem~5.3 cites [21], catalogue L123,
for earlier $\Sz(8)$ examples. That original is outside the fixed
corpus; no additional original is silently added. Our nine retained
matrix markings have separate literal panel certificates.

Theorem~5.2 quotes the earlier rank-two classification [18], catalogue
L115. The accepted common theory proves its normalizer families and
supplies the corrected subfield families. The precise row correspondence
is as follows, with $s$ denoting the parameter used in these tables:
\begin{center}\small
\begin{tabular}{p{26mm}p{109mm}}\toprule
Source rows & Treatment \\\midrule
Table 2, 1--4 & Split-prime normalizer families, intersection $C_2$.\\
Tables 2/3, 5 & Universal split-normalizer/affine pair, intersection $C_{s-1}$.\\
Tables 2/3, 6--11 & The two nonsplit normalizer families; opposite-factor and two-group partners are included.\\
Table 4, 12 & The same universal split-normalizer/affine pair.\\
Table 4, 13--15 & Full torus normalizer and maximal subfield, meeting in its compatible full normalizer at prime index.\\
Table 4, 18--20 & The corresponding torus extensions short of the full ambient normalizer.\\
Table 4, 17 & Distinct maximal simple subfields with a common simple subfield of parameter at least eight.\\
Table 4, 16 & The printed full central factor fails maximality; replace it by the complete scalar-plane family.\\\bottomrule
\end{tabular}
\end{center}
For row 16, with $K=\F_t\subset F=\F_s$ and $B_t$ the subfield
point stabilizer, the printed second parabolic is $Z_FB_t$.
For any $K$-plane $W$ properly contained in $F$,
\[
 B_t<Z_WB_t<Z_FB_t.
\]
Thus the printed pair is not residually weakly primitive. The complete
correct family consists of $Z_WB_t$ with $\dim_K W=2$; it is locally
doubly transitive. The accepted subfield theorem proves exhaustion and
its field-orbit equivalence, rather than just producing omitted examples.

The local column in Table~2 is correct in row 5 and in rows 9 and 11
at $s=8$; the printed $q=8$ in row 11 must refer to this table's $s$.
Table~3 retains row 5 as its only local case. A cyclic point stabilizer
of order four can be transitive on all other points of a prime torus
action only at degree five, which occurs for the minus torus at $s=8$.
The corrected central families are additional local cases in Table~4.

The preliminary note proves or explicitly qualifies each numbered
Lemma~2.1 and 3.1--3.11. Four literal errors must be retained:
Lemma~3.5 fails at $(q,s)=(32768,8)$; its usable conclusion concerns
a prime quotient, not the compositeness of $s-1$.
Lemma~3.6(ii) needs field-degree divisibility, not integer divisibility
of field sizes: $(q,s)=(32,8)$ disproves the printed version.
Lemma~3.6(iii) uses the residue of $e$ modulo four, not $2e+1$.
Lemma~3.9(ii) must keep the prime extension in the compatible torus;
$q=8$, $d=5$, $p=13$ disproves the opposite-factor allowance.
Lemma~3.6(i) reads $5\mid q^2+1$ on the native page; a previous OCR
reading of $s^2+1$ was our transcription issue. The overgroup statements
3.10 and 3.11 concern strict proper overgroups. Lemma~3.8 is explicitly
the source's cited Aschbacher theorem with its diagram hypotheses.

\section{Realizations, flag transitivity and the failed exhaustive list}
The separate realization note reproduces the earlier reviewed uniform
arguments, now with the uniqueness and scalar-field interfaces explicit.
Its conclusions for the individual source patterns are:
\begin{center}
\begin{tabular}{lcc}\toprule
Pattern & Intersections exist & RWPRI realization\\\midrule
Theorem 5.3, all three & yes & yes\\
Theorem 5.5, 1, 2, 3, 5 & no & no\\
Theorem 5.5, 4, 6, 7, 8, 10 & yes & yes\\
Theorem 5.5, 9 & yes & no\\\bottomrule
\end{tabular}
\end{center}
The mixed pattern 9 has forced full central coordinates; intermediate
scalar subspaces destroy maximality in every realization of that full
pattern. This is stronger than the failure of one convenient construction.
These arguments answer the closing realization question for each
printed pattern. They do not make the printed list exhaustive.

The existing compatible-tori example has parabolics
$\Sz(8)$, $\Sz(32)$ and $C_{533}:C_4$ in $\Sz(32768)$.
Their intersections have orders $20,52,164$, with final intersection
of order four. It is a thick, flag-transitive, residually connected
RWPRI geometry and violates the applicable branch of Theorem~5.5.
A permutation of types cannot turn its order-twenty intersection into
the printed order-eight one. This remains the already registered
principal counterexample.

The source's positive flag-transitivity remarks can be proved without
assuming that a subgroup pattern already has a geometric diagram.
If three parabolics $A,B,C$ satisfy
$A=(A\cap B)(A\cap C)$, every chamber has a common representative:
after translating its $A$-coset to $A$, write the other cosets as
$bB,cC$ with $b,c\in A$ and factor $b^{-1}c=xy$ with
$x\in A\cap B$, $y\in A\cap C$. Then $bx=cy^{-1}$ lies in all
three. Hence the geometry is flag transitive. It suffices to check
\[
 |A\cap B|\,|A\cap C|/|A\cap B\cap C|=|A|.
\]
This identity holds in the split dihedral parabolic of Theorem~5.3's
second pattern and in its order-four parabolic when the third pattern
has $p=2$. It also holds in the soluble third parabolic of the viable
composite patterns 4, 6, 7 and 8. The other composite patterns covered
by the source's remark are impossible, so make no positive existence
claim. Patterns 9 and 10 were expressly excepted there.

No viable printed composite pattern is algebraically locally doubly
transitive. In patterns 4 and 6 the cyclic prime quotient has odd
degree and a regular action. In pattern 7, the action of $\Sz(t)$
on its split-normalizer cosets has degree greater than the stabilizer
order plus one. In pattern 8 the two distinct extension primes would
both have to be five: each local action has a point stabilizer whose
image has order dividing four. In pattern 10 a simple subfield
$\Sz(u)$ cannot be a doubly transitive point stabilizer in
$\Sz(u^r)$ for $r\ge3$: using
$x^5/2<|\Sz(x)|<x^5$ shows its coset degree exceeds
$|\Sz(u)|+1$. These arguments concern only the displayed types.
They cannot exclude the omitted scalar-subspace families.

\section{Every clause of the local-transitivity conclusion}
Theorem~6.1 has distinct assertions, and each must be treated separately.
The retained universal rank-two pair exists and has the stated
properties. Its asserted uniqueness for composite $q-1$ is false:
the accepted scalar-plane family already gives valencies $65,8$ at
$q=32768$, distinct from the source pair's valencies $2,32768$.
The prime-$q-1$ rank bound and rank-two counts survive, as proved above.

Thinness fails in both prime branches under the literal algebraic
local definition. In $\Sz(8)$ an accepted marking has parabolic
orders $14,10,26$; in $\Sz(32)$ another has orders $62,62,82$.
Every actual panel has three elements, whereas the algebraic standard
orbit has two. Independent polynomial-matrix readers check the complete
small subgroup and coset lists, all intersections, and the explicit
extra representatives. Each is a rank-three geometry satisfying RWPRI
and algebraic local double transitivity. Thus neither the $q=8$ clause
nor the prime $q-1>7$ clause gives actual thinness or flag transitivity.
The nine earlier thin markings remain valid sufficient examples.

Finally the separate universal bound on locally doubly transitive rank
is false. The new note \emph{Scalar subspaces and unbounded algebraic
local rank} chooses independent prime subfields over $K=\F_t$ and
$T=K+Kx_1+\cdots+Kx_k$. Its parabolics are the $k$ maximal
subfields omitting one coordinate and the group $Z_TB_t$.
Every intersection deletes exactly the corresponding coordinates.
Scalar hyperplanes give maximal inclusions; the last local actions
are $\Sz(t)/B_t$ and the affine action on $K$. This proves the
source's two algebraic conditions in rank $k+1$ for arbitrary $k$.
For $k=2$ the coset system is a geometry with composite $q-1$;
for larger $k$ no unnecessary geometry or flag-transitivity assertion
is made. The source explicitly excludes pregeometries of these ranks,
so this is enough to refute its stated bound.

\section{Attribution, verification and limits}
Original numbered statements and examples are attributed to the fixed
source and its cited predecessors. Classical Suzuki structure remains
an explicit premise. The common scalar-subspace theory and earlier
standard-subfield and compatible-tori proofs are reused with their
accepted provenance. The new contributions in this completion are the
necessity arguments, complete source joins and corrections, two literal
thinness witnesses, and the independent-coordinate rank construction.
No priority claim is made for the elementary lemmas.

The proof volume preserves the exact accepted component texts, including
historical sentences that described obligations still open when a
component was written. The present ledger supersedes those progress
statements; it does not silently edit the frozen proofs. In particular
the prime-pattern note's then-unresolved thinness obligation is disposed
of by the two subsequent witnesses and the rank note. Our earlier
right-coset wording in the first thinness draft was a transcription
mistake, corrected before acceptance; the source already uses left cosets.

The independent certificate readers establish finite matrix and coset
identities. Assembly checks establish exact attribution, component
bindings and preservation of the native proof pages. Neither kind of
software check substitutes for the distinct mathematical review of the
general proofs. There is no new scientific search, Lean formalization,
compression measurement, claim of a complete corrected classification,
or recovery credit for an excluded cited original. A statement with
stronger actual-residue or flag-transitivity hypotheses is a different
problem and is not declared settled by the algebraic rank construction.
\end{document}
